\begin{Exercise}[label={ex:compressa}]
    Si considerino le seguenti definizioni:
    \begin{lstlisting}[style=cstylesmall]
        // Lista di interi
        typedef struct el{
            int dato;
            struct el* prox;
        } Node;
        typedef Node* Lista;

        // Lista compressa di interi
        typedef struct c_el{
            int dato, count;
            struct c_el* prox;
        } CNode;
        typedef CNode* ListaCompressa;
    \end{lstlisting}
    Le liste compresse rappresentano le classiche liste con nodi in cui appare il valore e il numero di volte che appare nella lista originale.
    
    \textbf{Esempio:} la lista \texttt{[3,3,3,5,5,2,2,2,2]} viene rappresentata come \texttt{[(3,3),(5,2),(2,4)]}.  
    Scrivere una funzione che prende una lista e costruisce la sua rappresentazione compressa.
\end{Exercise}

\begin{Answer}[ref={ex:compressa}]
  	\lstinputlisting[style=cstylesmall]{src/compressa.c}
\end{Answer}